% Rozdział 3 — Metodologia badań i projekt systemu
% Konwencje: numeracja sekcji z klasy dokumentu (plik zakłada \chapter/\section
% w szablonie pracy); istniejące pozycje bibliografii cytowane literalnie [n]
% jak w całej pracy; NOWE pozycje przez \cite{...} — klucze w bibliografia_nowe.tex.
% Nazwy cech/kolumn w \texttt{...} z eskejpowanym podkreślnikiem.

\section{Cel i zakres badań}

Celem pracy jest zbadanie, czy i w jakim stopniu dynamiczne, kalendarzowe
monitorowanie ryzyka niewykonania zobowiązania przewyższa klasyczną, jednorazową
ocenę statyczną. Rozdziały 1 i 2 wykazały, że zarówno wymogi regulacyjne (art.~70
Prawa bankowego, wytyczne EBA dotyczące monitorowania kredytów [9]), jak i rozwój
metod uczenia maszynowego przesuwają akcent z pojedynczej decyzji kredytowej ku
ciągłej obserwacji ekspozycji. Część badawcza operacjonalizuje tę obserwację
w~postaci schematu, w~którym ta sama ekspozycja oceniana jest wielokrotnie na
przesuwanym oknie historii płatniczej, a~wynikiem jest trajektoria
prawdopodobieństwa niewykonania zobowiązania (PD) zamiast pojedynczego wyniku.

Zakres badań obejmuje: (1)~konstrukcję panelu danych z~przesuwanym oknem na bazie
publicznego zbioru UCI \emph{default of credit card clients} [13];
(2)~trening i~kalibrację pięciu klasyfikatorów reprezentujących odmienne rodziny
algorytmiczne (las losowy, XGBoost, LightGBM, CatBoost --- ujęcie statyczne;
LSTM --- ujęcie sekwencyjne); (3)~budowę kompletnego systemu eksperymentalnego
(frontend, backend, serwis predykcyjny, baza danych) realizującego monitoring
trajektorii wraz z~progami alertu i~wyjaśnieniami predykcji; (4)~ewaluację
porównawczą reguły statycznej i~monitorującej oraz audyt sprawiedliwości.
Poza zakresem pozostają: dane dochodowe (nieobecne w~zbiorze --- ograniczenie
omówione we Wstępie), dane alternatywne oraz zespoły typu stacking (świadoma
decyzja zakresu; kierunek dalszych badań).

\section{Postawione hipotezy badawcze}

Pytania badawcze sformułowane we Wstępie skonkretyzowano w~postaci trzech
weryfikowalnych hipotez:

\textbf{H1 (zachowanie jakości przy krótszym oknie).} \emph{Modele trenowane na
3-miesięcznym oknie obserwacji (W3) zachowują zdolność dyskryminacyjną modeli
6-miesięcznych: strata AUC nie przekracza 1~punktu procentowego.} Hipoteza
warunkuje wykonalność całego schematu monitoringu --- okno 3-miesięczne jest ceną
za możliwość zbudowania czterech punktów trajektorii z~dostępnych 6~miesięcy
historii.

\textbf{H2 (wartość monitoringu).} \emph{Kalendarzowa reguła monitorująca ---
alert, gdy PD w~którymkolwiek z~okien W0..W3 przekracza próg --- oferuje
wcześniejszą detekcję pogarszającej się sytuacji dłużnika oraz wykrycia
niedostępne regule statycznej (opartej wyłącznie na najnowszym oknie W3), przy
porównywalnej ogólnej zdolności dyskryminacyjnej.} Hipoteza celowo nie postuluje
wyższej czułości przy stałym budżecie fałszywych alarmów: agregacja maksimum po
czterech skorelowanych oknach zwiększa ekspozycję na szum, a~wartością
monitoringu jest wyprzedzenie czasowe i~komplementarność, nie substytucja oceny
statycznej.

\textbf{H3 (sprawiedliwość).} \emph{Modele zachowują parytet względem atrybutu
chronionego SEX: różnica parytetu demograficznego (DPD) oraz różnica wyrównanych
szans (EOD) nie przekraczają 0,10 przy binaryzacji progami produkcyjnymi.}
Wybór metryk i~progu omówiono w~rozdziale~5.5; wymóg nawiązuje do klasyfikacji
scoringu kredytowego jako zastosowania wysokiego ryzyka w~rozporządzeniu
o~sztucznej inteligencji [25].

Weryfikację hipotez przedstawia rozdział~5.6.

\section{Struktura i charakterystyka danych dłużników}

\subsection{Źródła i sposób pozyskania danych}

Wykorzystano publiczny zbiór \emph{default of credit card clients} (UCI Machine
Learning Repository), obejmujący 30~000 posiadaczy kart kredytowych tajwańskiego
banku [13]. Każdy rekord zawiera 23~zmienne objaśniające oraz binarną etykietę
niewykonania zobowiązania w~październiku 2005~r. (22,1\% obserwacji pozytywnych).
Zbiór jest standardowym benchmarkiem w~literaturze credit scoringu, co umożliwia
odniesienie wyników do wcześniejszych badań [10], [13].

\subsection{Opis zmiennych wejściowych i zmiennej docelowej}

Zmienne dzielą się na trzy grupy: (1)~profil klienta --- limit kredytowy
(\texttt{LIMIT\_BAL}), płeć (\texttt{SEX}: 1~=~mężczyzna, 2~=~kobieta),
wykształcenie (\texttt{EDUCATION}), stan cywilny (\texttt{MARRIAGE}), wiek
(\texttt{AGE}); (2)~historia płatnicza za sześć kolejnych miesięcy
(kwiecień--wrzesień 2005): statusy płatności \texttt{PAY\_*}, salda rachunków
\texttt{BILL\_AMT1..6}, kwoty wpłat \texttt{PAY\_AMT1..6}; (3)~zmienna docelowa
\texttt{Default}. Istotną cechą zbioru jest przesunięta numeracja: najnowszy
miesiąc (wrzesień) opisują kolumny \texttt{PAY\_0}, \texttt{BILL\_AMT1}
i~\texttt{PAY\_AMT1}, a~kolumna \texttt{PAY\_1} nie istnieje. Mapowanie kolumn
na oś czasu przedstawia tabela~\ref{tab:mapowanie-kolumn}.

\begin{table}[htbp]
\centering
\caption{Mapowanie kolumn UCI na miesiące kalendarzowe.}
\label{tab:mapowanie-kolumn}
\begin{tabular}{llll}
\toprule
Miesiąc & Status płatności & Saldo rachunku & Kwota wpłaty \\
\midrule
wrzesień (najnowszy) & \texttt{PAY\_0} & \texttt{BILL\_AMT1} & \texttt{PAY\_AMT1} \\
sierpień             & \texttt{PAY\_2} & \texttt{BILL\_AMT2} & \texttt{PAY\_AMT2} \\
lipiec               & \texttt{PAY\_3} & \texttt{BILL\_AMT3} & \texttt{PAY\_AMT3} \\
czerwiec             & \texttt{PAY\_4} & \texttt{BILL\_AMT4} & \texttt{PAY\_AMT4} \\
maj                  & \texttt{PAY\_5} & \texttt{BILL\_AMT5} & \texttt{PAY\_AMT5} \\
kwiecień (najstarszy)& \texttt{PAY\_6} & \texttt{BILL\_AMT6} & \texttt{PAY\_AMT6} \\
\bottomrule
\end{tabular}
\end{table}

\subsection{Preprocessing danych: czyszczenie, normalizacja, inżynieria cech}

Dla modeli statycznych z~każdego okna obserwacji wyliczane jest 13~cech
pochodnych (m.in. \texttt{PAY\_mean}, \texttt{PAY\_max}, \texttt{BILL\_mean},
\texttt{BILL\_std}, \texttt{BILL\_trend}, \texttt{utilization\_rate} =
\texttt{BILL\_mean}~/~\texttt{LIMIT\_BAL}, \texttt{payment\_ratio},
\texttt{late\_count}, \texttt{severe\_late}, \texttt{recent\_pay\_status}),
uzupełnionych o~9~surowych kolumn okna oraz zmienne demograficzne kodowane
one-hot. Funkcja \texttt{engineer\_features(df, window)} jest sparametryzowana
oknem --- te same definicje cech obowiązują dla każdego z~okien W0..W3,
a~identyczna kopia modułu działa w~treningu i~w~serwisie inferencyjnym (spójność
trening--serwing jest przypięta testami automatycznymi). Kodowanie one-hot używa
stałych dziedzin kategorii zbioru UCI, dzięki czemu zestaw kolumn nie zależy od
składu przetwarzanej próbki --- w~szczególności od pojedynczego rekordu
w~inferencji. Wartości nieskończone i~brakujące powstające w~cechach ilorazowych
(np. przy \texttt{LIMIT\_BAL}~=~0) są zerowane identycznie po obu stronach.
Standaryzacja (StandardScaler) jest dopasowywana wyłącznie na części treningowej
i~stosowana przez transformację do pozostałych podzbiorów oraz w~inferencji.
Model sekwencyjny otrzymuje surowy tensor $3\times3$ (trzy miesiące $\times$
trzy kanały: PAY, BILL\_AMT, PAY\_AMT), skalowany per kanał skalerami
dopasowanymi na części treningowej --- bez cech pochodnych i~bez zmiennych
demograficznych.

\subsection{Konstrukcja panelu z przesuwanym oknem (W0..W3)}

Rdzeniem części badawczej jest przekształcenie przekrojowego zbioru UCI w~panel
umożliwiający symulację monitoringu. Z~sześciu miesięcy historii budowane są
cztery nakładające się okna 3-miesięczne, tworzące pseudo-oś czasu
(tabela~\ref{tab:okna}).

\begin{table}[htbp]
\centering
\caption{Definicje okien przesuwnych.}
\label{tab:okna}
\begin{tabular}{lllll}
\toprule
Okno & Zakres & Statusy & Salda & Wpłaty \\
\midrule
W0 (najstarsze) & kwi--cze & \texttt{PAY\_6,5,4} & \texttt{BILL\_AMT6,5,4} & \texttt{PAY\_AMT6,5,4} \\
W1              & maj--lip & \texttt{PAY\_5,4,3} & \texttt{BILL\_AMT5,4,3} & \texttt{PAY\_AMT5,4,3} \\
W2              & cze--sie & \texttt{PAY\_4,3,2} & \texttt{BILL\_AMT4,3,2} & \texttt{PAY\_AMT4,3,2} \\
W3 (najnowsze)  & lip--wrz & \texttt{PAY\_3,2,0} & \texttt{BILL\_AMT3,2,1} & \texttt{PAY\_AMT3,2,1} \\
\bottomrule
\end{tabular}
\end{table}

Obowiązuje zasada nadrzędna: \textbf{żadna wartość nie jest fabrykowana}. Każde
okno to rzeczywisty, trzymiesięczny wycinek historii klienta; jedyną operacją
jest wybór miesięcy ,,widocznych'' w~danym punkcie pseudo-czasu. Klient ma jedną
prawdziwą etykietę (default w~październiku), wspólną dla wszystkich okien ---
nie są tworzone żadne etykiety pośrednie.

Modele trenowane są wyłącznie na oknie W3, wyrównanym z~etykietą (,,mając
ostatnie 3~miesiące, czy klient nie wykona zobowiązania w~następnym
miesiącu?''). W~inferencji ten sam model stosowany jest do każdego z~okien
W0..W3 --- ponieważ każde okno ma identyczną strukturę 3-miesięcznego wycinka,
rozkład wejść inferencyjnych odpowiada treningowemu. Trening na jednym oknie na
klienta eliminuje ryzyko wycieku etykiet, które niosłaby augmentacja wszystkimi
czterema oknami.

\textbf{Ograniczenie metodologiczne (istotne dla interpretacji wyników).} Zbiór
UCI jest przekrojowy, więc monitoring jest symulowany retrospektywnie: cztery
okna tego samego klienta nakładają się w~dwóch trzecich, a~wszystkie przewidują
tę samą październikową etykietę --- okno W0 odpowiada predykcji
z~czteromiesięcznym horyzontem, W3 z~jednomiesięcznym. Trajektoria PD stanowi
zatem przybliżenie scenariusza produkcyjnego, a~nie jego pełną replikę;
walidacja na prawdziwym panelu podłużnym pozostaje kierunkiem dalszych badań.
Diagnozę odróżniającą narastanie ryzyka od ewentualnego przesunięcia rozkładu
między oknami przedstawia rozdział~5.4.

\subsection{Podział zbioru i protokół walidacji}

Zastosowano trójdzielny, stratyfikowany podział 60/20/20
(\texttt{random\_state}~=~42): 18~000 obserwacji treningowych, 6~000
kalibracyjnych i~6~000 testowych. Część kalibracyjna służy wyłącznie dopasowaniu
kalibracji izotonicznej (rozdz.~4.6) oraz wyznaczeniu progów alertu
(rozdz.~4.7) --- dzięki temu zbiór testowy pozostaje zamrożony: nie uczestniczy
w~treningu, doborze hiperparametrów, kalibracji ani doborze progów, a~wchodzi do
analizy wyłącznie w~finalnej ewaluacji rozdziału~5. Standaryzacja cech
dopasowywana jest na części treningowej (por.~3.3.3). Stratyfikacja zachowuje
w~każdym podzbiorze proporcję klas 78/22.

\begin{figure}[htbp]
\centering
\includegraphics[width=.95\textwidth]{fig_4_1_podzial_danych}
\caption{Podział danych: trening / kalibracja / test (60/20/20). Źródło: opracowanie własne.}
\label{fig:podzial}
\end{figure}

\section{Projekt architektury systemu eksperymentalnego}

\subsection{Architektura ogólna i przepływ danych}

System składa się z~czterech warstw komunikujących się przez HTTP/JSON
i~uruchamianych w~kontenerach (docker-compose; frontend w~trybie deweloperskim
poza kompozycją):

\begin{verbatim}
React 19 + TypeScript (5173)
   |  POST/GET /api/v1/monitoring/*
   v
ASP.NET Core .NET 8 (5120) - MonitoringController -> MonitoringService
   |  walidacja 22 cech, mapowanie nazw, labelki okien, obsluga bledow,
   |  persystencja (EF Core, transakcja atomowa)
   +--> Flask (5001): POST /predict/timeseries - 5 modeli x 4 okna,
   |        trendy, progi kosztowe, SHAP top-5
   +--> PostgreSQL 16 (5432): Client / Snapshot / Prediction / Trend
\end{verbatim}

Przepływ oceny monitorującej: frontend przesyła 22~cechy klienta wraz z~datą
migawki; backend waliduje żądanie, wywołuje serwis predykcyjny, wzbogaca
odpowiedź o~etykiety kalendarzowe okien i~--- w~ścieżce stanowej --- zapisuje
migawkę, predykcje okna W3 oraz trendy w~bazie. Kontrakt API (cztery endpointy,
typy, kody błędów) został sformalizowany przed implementacją, co umożliwiło
równoległą pracę nad backendem i~frontendem przeciwko wspólnemu mockowi.

\subsection{Warstwa backendowa}

Backend pełni rolę bramy i~warstwy trwałości. Waliduje wejście (adnotacje
zakresów dla 22~cech), mapuje konwencje nazewnicze, wyznacza etykiety
kalendarzowe okien względem daty migawki i~mapuje błędy na ustrukturyzowaną
kopertę (\texttt{VALIDATION\_FAILED}, \texttt{CONFLICT} --- duplikat pary
klient/data, \texttt{CLIENT\_NOT\_FOUND},
\texttt{ML\_SERVICE\_ERROR}/\texttt{UNAVAILABLE}, \texttt{INTERNAL\_ERROR}).
Warstwa trwałości (EF~Core, Npgsql, automatyczne migracje przy starcie)
przechowuje encje: \texttt{Client} (unikalny identyfikator biznesowy),
\texttt{Snapshot} (data + pełne 22~cechy wejściowe), \texttt{Prediction}
(PD per model dla okna etykietowanego W3) i~\texttt{Trend} (nachylenie
trajektorii i~alert per model, aktualizowane w~trybie upsert). Zapis migawki
wraz z~pięcioma predykcjami i~pięcioma trendami wykonywany jest w~jawnej
transakcji bazodanowej --- awaria w~trakcie zapisu nie może pozostawić migawki
bez predykcji; właściwość ta jest weryfikowana testem integracyjnym
z~wstrzykniętą awarią na rzeczywistym PostgreSQL.

\subsection{Integracja modeli AI z systemem}

Serwis predykcyjny (Flask) ładuje przy starcie pięć skalibrowanych artefaktów
W3 (cztery modele drzewiaste \texttt{.pkl} --- opakowania
\texttt{CalibratedClassifierCV} --- oraz model LSTM \texttt{.keras}
z~zewnętrznym kalibratorem izotonicznym), skalery i~plik progów alertu.
Endpoint \texttt{POST /predict/timeseries} przyjmuje 22~cechy, wewnętrznie
mapuje wartości każdego z~okien W0..W3 w~sloty kolumn W3 (dzięki czemu model
trenowany na W3 ocenia każde okno w~tej samej przestrzeni cech), i~zwraca:
trajektorię (5~modeli $\times$ 4~okna), trendy z~regułą alertu opartą na
nachyleniu (W3~$-$~W0), progi kosztowe, flagi alertów per okno oraz wyjaśnienia
SHAP (pięć najistotniejszych cech per model drzewiasty, dla okna W3).
Historyczny endpoint \texttt{POST /predict} (modele 6-miesięczne bez
kalibracji) zachowano jako punkt odniesienia --- oba tryby są rozdzielone
i~nie mieszają artefaktów.

\subsection{Warstwa frontendowa}

Interfejs (React~19, TypeScript, Vite; wykresy Recharts) ma dwie zakładki.
Zakładka \emph{Prediction} odtwarza klasyczną, jednorazową ocenę (formularz
22~cech, karty wyników per model). Zakładka \emph{Monitoring} realizuje
Wariant~B: lista monitorowanych klientów z~rozwijaną historią (master--detail),
wykres trajektorii PD (pięć linii, cztery okna), karty alertów semaforowych per
model, formularz datowanej migawki z~funkcją kopiowania poprzednich wartości
oraz komponent wyjaśnień SHAP (poziome słupki rozbieżne: wkład dodatni podnosi
PD). Warstwa typów TypeScript odwzorowuje kontrakt API w~pełnej,
pięciomodelowej postaci.

\section{Obsługa wyjątków i scenariusze brzegowe}

Obsługę błędów zaprojektowano w~trzech warstwach. Serwis ML odrzuca żądania
z~brakującym lub nienumerycznym polem (kod \texttt{VALIDATION\_FAILED}) ---
walidacja typów zapobiega zarówno awarii ścieżki sekwencyjnej, jak i~cichej
imputacji zera w~ścieżce statycznej (zero jest w~tym zbiorze znaczącą wartością
statusu płatności). Backend mapuje niedostępność serwisu ML na 503, jego błędy
na 502, duplikat migawki na 409, nieznanego klienta na 404, a~scenariusze
brzegowe cech ilorazowych (np. zerowy limit) neutralizuje identycznie jak
trening. Frontend tłumaczy kody błędów na komunikaty użytkownika. Scenariusze
te są pokryte testami integracyjnymi wszystkich trzech warstw.

\section{Narzędzia i środowisko technologiczne}

Stack: Python~3.11 (scikit-learn, XGBoost, LightGBM, CatBoost,
TensorFlow/Keras, fairlearn, SHAP, Optuna), Flask; .NET~8 (ASP.NET Core,
EF~Core, Npgsql); React~19 + TypeScript + Vite; PostgreSQL~16; Docker Compose.
Testy: pytest (serwis ML --- w~tym testy parytetu trening/serwing), xUnit
(backend --- w~tym testy integracyjne na rzeczywistym PostgreSQL przez
Testcontainers), Vitest (frontend). Metodyka pracy: GitHub Flow --- gałęzie
funkcjonalne, przegląd drugiego autora dla każdego pull requesta, ciągła
integracja (cztery zadania: backend, serwis ML, warstwa treningowa, frontend)
blokująca scalenie przy czerwonych testach. Zakres prac zdekomponowano na
zadania z~jawnym grafem zależności i~definicją ukończenia, prowadzone
w~sześciu dwutygodniowych sprintach planistycznych. Wszystkie wyniki liczbowe
rozdziału~5 są reprodukowalne pojedynczym uruchomieniem skryptów warstwy
treningowej, a~trening modelu sekwencyjnego jest deterministyczny (ustalone
ziarno losowe).
